% SD4H Appendix — unlimited pages
% \cmark and \xmark defined in main.tex preamble (via pifont)

% ============================================================
% APPENDIX A: Dataset Statistics
% ============================================================
\section{Dataset Statistics}
\label{app:datasets}

\begin{table}[h]
\centering
\caption{ICU databases used in this study. All data is preprocessed via the YAIB benchmark~\citep{vandewater2024yaib} with standardized cohort definitions, feature extraction, and train/validation/test splits.}
\label{tab:datasets}
\footnotesize
\resizebox{\columnwidth}{!}{%
\setlength{\tabcolsep}{3pt}
\begin{tabular}{@{}llrrll@{}}
\toprule
\textbf{Dataset} & \textbf{Role} & \textbf{Stays} & \textbf{Hosp.} & \textbf{Country} & \textbf{Source} \\
\midrule
\eicu{}   & Source   & 200,859 & 208 & US & Telehealth net. \\
\mimic{}  & Target   & 73,181  & 1   & US & Academic med.\ ctr. \\
\hirid{}  & Source 2 & 33,905  & 1   & Switz. & Univ.\ hospital \\
\bottomrule
\end{tabular}%
}
\end{table}

\begin{table*}[h]
\centering
\caption{Clinical prediction tasks. MI = missing indicator features (binary flags for whether a value was measured or forward-filled). Per-timestep tasks produce a prediction at every hourly timestep; per-stay tasks produce a single prediction per ICU stay.}
\label{tab:tasks}
\small
\begin{tabular}{@{}llllrl@{}}
\toprule
\textbf{Task} & \textbf{Type} & \textbf{Label} & \textbf{Granularity} & \textbf{Feature Dim} & \textbf{Pos.\ Rate} \\
\midrule
Mortality24 & Classification & ICU death within 24\,h & Per-stay & 96 (48 + 48 MI) & $\sim$5.5\% \\
AKI & Classification & KDIGO Stage $\geq$1 & Per-timestep & 100 (48 + 48 MI + 4 static) & $\sim$12\% \\
Sepsis & Classification & Sepsis-3 criteria & Per-timestep & 100 (48 + 48 MI + 4 static) & $\sim$1.1\% \\
LoS & Regression & Hours remaining & Per-timestep & 52 (48 + 4 static) & N/A \\
KF & Regression & Serum creatinine (mg/dL) & Per-stay & 292 (48 + 48 MI + 192 gen + 4 static) & N/A \\
\bottomrule
\end{tabular}
\end{table*}

The 48 clinical features include vital signs (heart rate, blood pressure, SpO$_2$, temperature, respiratory rate) and laboratory values (metabolic panels, blood gases, liver enzymes, cardiac markers, coagulation studies).
The 4 static features encode patient demographics (age, sex, height, weight).
KF uses 192 additional generated features: cumulative minimum, maximum, mean, and count statistics for each of the 48 clinical features, which must maintain internal consistency after adaptation.
LoS omits missing indicators because the YAIB LoS LSTM was trained on a 52-feature schema.


% ============================================================
% APPENDIX B: Implementation Details
% ============================================================
\section{Implementation Details}
\label{app:implementation}

\begin{table}[h]
\centering
\caption{Hyperparameters for the \retr{} architecture.}
\label{tab:hyperparams}
\footnotesize
\setlength{\tabcolsep}{3pt}
\begin{tabular}{@{}lp{0.55\columnwidth}@{}}
\toprule
\textbf{Hyperparameter} & \textbf{Value} \\
\midrule
$d_\text{model}$ & 128 \\
$d_\text{latent}$ & 128 \\
$n_\text{heads}$ & 8 \\
$d_\text{ff}$ & 512 \\
$n_\text{enc\_layers}$ & 4 \\
$n_\text{dec\_layers}$ & 3 \\
$n_\text{cross\_layers}$ & 2 (Mort., Sepsis) / 3 (AKI, LoS, KF) \\
$k$ (neighbors) & 16 \\
$W$ (retrieval window) & 6 timesteps (matches 6-hour clinical reassessment cycle) \\
Learning rate & $5 \times 10^{-4}$ \\
Optimizer & AdamW \\
Batch size & 16 \\
Phase~1 pretrain epochs & 10--50 (task-dependent) \\
Phase~2 training epochs & 30--100 (early stopping) \\
$\lambda_\text{fidelity}$ & 1.0 \\
$\lambda_\text{range}$ & 0.1 \\
Memory refresh ($R$) & 5 epochs \\
Dropout & 0.2 \\
\bottomrule
\end{tabular}
\end{table}

\paragraph{Task-specific settings.}
Temporal attention mode is set to \emph{causal} for per-timestep tasks (AKI, Sepsis, LoS) to prevent information leakage from future timesteps, and \emph{bidirectional} for per-stay tasks (Mortality, KF) where the full sequence is available at prediction time.
Cross-domain normalization (an affine transform matching per-feature mean and variance to target statistics, computed once from training data) is applied to all tasks.
Phase~1 self-retrieval is optionally enabled: during autoencoder pretraining, a memory bank is built from the encoder's own target representations and provided to cross-attention blocks, avoiding degenerate zero-context initialization. We enable self-retrieval for all tasks in this work.
Variable-length batching is used for per-timestep tasks and disabled for per-stay tasks.

\paragraph{Hardware.}
Primary development hardware: NVIDIA V100S-32GB.
Multi-seed validation also used A6000-48GB, RTX 3090-24GB, L40S-48GB, and A100-40GB (SLURM cluster).
Cross-server variance is within $\pm$0.003 \aucroc{} for all tasks.

\paragraph{Model sizes.}
The retrieval adapter contains 2.59M parameters (Mortality, Sepsis; $n_\text{cross}=2$) to 2.88M parameters (KF; $n_\text{cross}=3$).
Frozen LSTM baselines range from 170K (Sepsis) to 1.42M (AKI) parameters.
DA baselines: DANN 350K, Deep CORAL 251K, CoDATS 428K trainable parameters.

\paragraph{AdaTime configuration.}
The same \retr{} architecture is used with reduced capacity: $d_\text{model}=64$, $d_\text{latent}=64$, $d_\text{ff}=128$, $n_\text{enc\_layers}=2$, $n_\text{dec\_layers}=1$, $n_\text{cross\_layers}=2$--$3$.
The frozen backbone is AdaTime's standard 3-block 1D-CNN ($\sim$200K parameters).
Adapter size: 100K--450K parameters depending on dataset.
All AdaTime runs (ours and every baseline) follow the benchmark's exact recipe: 40-epoch budget, Adam with $\beta=(0.5, 0.99)$, weight decay $10^{-4}$, last-epoch model, no source-validation split.

\paragraph{AdaTime end-to-end baselines.}
\citet{ragab2023adatime} benchmark eleven state-of-the-art end-to-end DA methods: DANN, Deep CORAL, CDAN, DSAN, HoMM, MMDA, CoDATS, AdvSKM, DIRT-T, CoTMix, and SASA. We also run CLUDA~\citep{ozyurt2023cluda} and RAINCOAT~\citep{he2023raincoat} under the same protocol. For each dataset, \Cref{tab:adatime} reports the best of these published methods.

\paragraph{End-to-end DA on clinical tasks.}
As a stress test for the frozen-predictor assumption, we also ran three end-to-end DA methods that retrain the full LSTM on each clinical task: CLUDA~\citep{ozyurt2023cluda}, RAINCOAT~\citep{he2023raincoat}, and ACON~\citep{liu2024acon}. All three were swept over published hyperparameters and trained under right-padded YAIB data (matching the frozen-baseline protocol; left-padding inflated per-timestep AUROC and was excluded). \Cref{tab:e2e-clinical} reports the best hyperparameter per method per task. \retr{}, despite never touching the predictor, beats the best end-to-end DA on all three tasks.

\begin{table}[h]
\centering
\caption{End-to-end DA baselines on clinical tasks (right-padded, best hyperparameter per method). AUROC ($\times 100$, $\uparrow$). $\Delta$ vs.\ the frozen source-only baseline.}
\label{tab:e2e-clinical}
\footnotesize
\begin{tabular}{@{}l ccc@{}}
\toprule
\textbf{Method} & \textbf{Mortality} & \textbf{AKI} & \textbf{Sepsis} \\
\midrule
Frozen source-only            & 80.79 & 85.58 & 71.59 \\
\midrule
E2E RAINCOAT                   & 81.52 & 89.55 & 76.24 \\
E2E CLUDA                      & 81.81 & 87.35 & 76.61 \\
E2E ACON                       & 82.23 & 70.79 & 75.30 \\
\midrule
\retr{} (ours, frozen)         & \textbf{85.74} & \textbf{91.14} & \textbf{77.76} \\
\bottomrule
\end{tabular}
\end{table}


% ============================================================
% APPENDIX C: Extended Results
% ============================================================
\section{Extended Results}
\label{app:extended}

\subsection{AUCPR Results (eICU $\to$ MIMIC)}
\label{app:aucpr}

\begin{table}[h]
\centering
\caption{\aucpr{} results ($\times 100$) for classification tasks. Bootstrap 95\% CIs (500 replicates).}
\label{tab:aucpr}
\small
\resizebox{\columnwidth}{!}{%
\begin{tabular}{@{}lccc@{}}
\toprule
& \textbf{Frozen Baseline} & \textbf{\retr{}} [95\% CI] & $\boldsymbol{\Delta}$ \\
\midrule
Mortality & 29.67 & 34.88 [32.45, 37.36] & +5.21 \\
AKI       & 56.78 & 72.86 [72.65, 73.06] & +16.08 \\
Sepsis    &  2.98 &  5.32 [5.13, 5.51]   & +2.34 \\
\bottomrule
\end{tabular}%
}
\end{table}

AKI \aucpr{} gains (+16.08) are nearly 3$\times$ the \aucroc{} gains (+5.56): adaptation disproportionately improves precision at clinically relevant operating points.

\subsection{Regression Extended Results (eICU $\to$ MIMIC)}
\label{app:regression}

\begin{table}[h]
\centering
\caption{Regression results across multiple metrics. Bootstrap 95\% CIs (500 replicates). MAE and RMSE: lower is better ($\downarrow$); $R^2$: higher is better ($\uparrow$).}
\label{tab:regression}
\resizebox{\columnwidth}{!}{%
\footnotesize
\setlength{\tabcolsep}{3pt}
\begin{tabular}{@{}llccr@{}}
\toprule
\textbf{Task} & \textbf{Metric} & \textbf{Baseline} & \textbf{\retr{}} [95\% CI] & $\boldsymbol{\Delta}$ \\
\midrule
LoS & MAE ($\downarrow$) & 0.2527 & 0.2331 [.2329, .2333] & $-$0.0196 \\
LoS & RMSE ($\downarrow$) & 0.3123 & 0.3004 [.3001, .3007] & $-$0.0119 \\
LoS & $R^2$ ($\uparrow$) & 0.1130 & 0.1794 [.1782, .1806] & +0.0664 \\
\addlinespace
KF  & MAE ($\downarrow$) & 0.0330 & 0.0256 [.0249, .0262] & $-$0.0074 \\
KF  & RMSE ($\downarrow$) & 0.0600 & 0.0477 [.0451, .0503] & $-$0.0123 \\
KF  & $R^2$ ($\uparrow$) & 0.7489 & 0.8411 [.8266, .8553] & +0.0922 \\
\bottomrule
\end{tabular}%
}
\end{table}

Table~\ref{tab:regression} reports the primary retrieval variant for each task. The best task-specific variants (starred in Table~\ref{tab:main-results}) go further: LoS without MMD alignment achieves MAE 0.2212 (approximately 37.2 hours when denormalized), a 5.3-hour reduction from the frozen baseline's 42.5 hours and 2.0 hours below the eICU-native LSTM (39.2h); KF without fidelity and without MMD achieves normalized MAE 0.0239 ($\Delta = -0.0091$, 0.292 mg/dL denormalized).
KF $R^2$ improves from 0.749 to 0.841, nearly doubling the explained variance in serum creatinine predictions.
Dropping the fidelity loss removes fidelity gradient dominance in KF's 292-feature cumulative-statistics regime.

\subsection{HiRID $\to$ MIMIC Extended Results}
\label{app:hirid}

\begin{table}[h]
\centering
\caption{\hirid{} $\to$ \mimic{} results. Bootstrap 95\% CIs (500 replicates). All $p < 0.001$ except KF ($p = 0.482$, n.s.).}
\label{tab:hirid}
\small
\begin{tabular}{@{}llcl@{}}
\toprule
\textbf{Task} & \textbf{Metric} & $\boldsymbol{\Delta}$ \textbf{vs Frozen} [95\% CI] & $\boldsymbol{p}$ \\
\midrule
Mortality & \aucroc{} & +0.0488 [+0.0284, +0.0711] & $<$0.001 \\
Mortality & \aucpr{}  & +0.0502 [+0.0167, +0.0859] & 0.002 \\
\addlinespace
AKI & \aucroc{} & +0.0776 [+0.0745, +0.0807] & $<$0.001 \\
AKI & \aucpr{}  & +0.1471 [+0.1414, +0.1531] & $<$0.001 \\
\addlinespace
Sepsis & \aucroc{} & +0.0777 [+0.0707, +0.0841] & $<$0.001 \\
Sepsis & \aucpr{}  & +0.0525 [+0.0471, +0.0584] & $<$0.001 \\
\addlinespace
LoS & MAE & $-$0.0452 [$-$0.0497, $-$0.0410] & $<$0.001 \\
\addlinespace
KF & MAE & +0.0000 [$-$0.0009, +0.0010] & 0.482 \\
\bottomrule
\end{tabular}
\end{table}

\hirid{} results show that the adapter works across source domains with very different data collection practices (Swiss vs.\ US ICUs).
AKI and Sepsis improvements from \hirid{} (+0.078 \aucroc{} each) exceed those from \eicu{} (+0.056 and +0.062), possibly because greater domain distance provides more signal for adaptation.
KF is the only non-significant result, likely due to \hirid{}'s smaller cohort size (33,905 vs.\ 200,859 stays) providing insufficient coverage of the 292-feature KF schema.

\subsection{Seed Variance}
\label{app:seed}

\begin{table}[h]
\centering
\caption{Multi-seed \aucroc{} variance ($\times 100$). Spread = max $-$ min across seeds.}
\label{tab:seed}
\small
\begin{tabular}{@{}llrcr@{}}
\toprule
\textbf{Task} & \textbf{Config} & \textbf{Seeds} & \textbf{Mean $\Delta$} & \textbf{Spread} \\
\midrule
AKI & V5, $n_\text{cross}=3$ & 5 & +5.47 & 0.27 \\
Mortality & V4+MMD & 3 & +4.56 & 0.29 \\
Sepsis & V4+MMD & 3 & +4.94 & 2.03 \\
\bottomrule
\end{tabular}
\end{table}

AKI and Mortality show low variance (spread $<$0.3 \aucroc{} points), both well within noise.
AKI's 5-seed validation spans 3 physically distinct GPU servers (V100S, A6000, RTX 3090), with cross-server variance adding only 0.2 points.
Sepsis exhibits higher variance (spread 2.03) due to its 1.1\% label density; this is inherent to the task rather than method-specific, as DA baselines show comparable variance.


% ============================================================
% APPENDIX D: Architecture-Agnostic Transfer
% ============================================================
\section{Architecture-Agnostic Transfer}
\label{app:mas}

To test whether the adapter learns a domain-invariant input mapping rather than LSTM-specific artifacts, we apply an LSTM-trained adapter \emph{zero-shot} to frozen GRU and TCN architectures (no retraining).

\begin{table}[h]
\centering
\caption{Architecture-agnostic transfer: LSTM-trained adapter applied zero-shot to frozen GRU and TCN predictors. \aucroc{} improvement ($\Delta \times 100$) over each architecture's frozen baseline. Percentage indicates fraction of LSTM gain retained.}
\label{tab:mas}
\footnotesize
\setlength{\tabcolsep}{4pt}
\begin{tabular}{@{}lccc@{}}
\toprule
\textbf{Task} & \textbf{LSTM $\Delta$} & \textbf{GRU $\Delta$ (\%)} & \textbf{TCN $\Delta$ (\%)} \\
\midrule
Mortality & +4.76 & +4.04 (85\%) & +3.42 (72\%) \\
AKI       & +5.56 & +3.11 (56\%) & +3.16 (57\%) \\
Sepsis    & +5.12 & +2.40 (47\%) & +4.42 (86\%) \\
\bottomrule
\end{tabular}
\end{table}

All six source--architecture combinations show positive transfer, with gains ranging from 47\% to 86\% of the LSTM improvement.
The strongest retention occurs for TCN on Sepsis (86\%) and GRU on Mortality (85\%); the input-space corrections appear largely architecture-independent.
This supports the interpretation that the adapter resolves distributional mismatches (e.g., forward-fill imputation artifacts) at the input level, producing features that benefit any predictor trained on target-domain data.


% ============================================================
% APPENDIX E: Ablation Study
% ============================================================
\section{Ablation Study}
\label{app:ablation}

We ablate the main components of the retrieval adapter.
Each row modifies exactly one setting from the control configuration.

\begin{table}[h]
\centering
\caption{Component ablation on AKI (\eicu{} $\to$ \mimic{}, \aucroc{} $\times 100$). Each row removes exactly one component from the full retrieval adapter ($n_\text{cross}=3$).}
\label{tab:ablation-aki}
\small
\begin{tabular}{@{}lrr@{}}
\toprule
\textbf{Configuration} & \textbf{$\Delta$ \aucroc{}} & \textbf{Effect} \\
\midrule
Full method (control) & +5.48 & --- \\
No retrieval (remove memory bank) & +4.88 & $-$0.60 \\
No feature gate & +5.02 & $-$0.46 \\
No MMD alignment & +5.19 & $-$0.29 \\
No fidelity loss & \textbf{+5.70} & +0.22 \\
No target normalization & +5.43 & $-$0.05 \\
\bottomrule
\end{tabular}
\end{table}

\begin{table}[h]
\centering
\caption{Component ablation on Sepsis (\eicu{} $\to$ \mimic{}, \aucroc{} $\times 100$). Control: full retrieval adapter with $n_\text{cross}=2$.}
\label{tab:ablation-sepsis}
\small
\begin{tabular}{@{}lrr@{}}
\toprule
\textbf{Configuration} & \textbf{$\Delta$ \aucroc{}} & \textbf{Effect} \\
\midrule
Full method (control) & +4.99 & --- \\
No smoothness loss & +4.75 & $-$0.24 \\
Oversampling (20$\times$) & $-$3.31 & $-$8.30 \\
$n_\text{cross}=3$ (vs 2) & +4.35 & $-$0.64 \\
\bottomrule
\end{tabular}
\end{table}

\paragraph{Key findings.}
\emph{Pretraining is essential}: removing Phase~1 (the autoencoder pretrain) collapses every classification task, with Mortality falling $10$ \aucroc{} points below the unadapted frozen baseline.
\emph{Retrieval provides consistent benefit}: removing the memory bank reduces AKI by $0.60$ \aucroc{} and all other tasks show similar degradation.
\emph{Oversampling is harmful for sparse labels}: 20$\times$ oversampling of positive examples causes catastrophic performance loss on Sepsis, likely due to gradient bias from repeated rare-event samples.


% ============================================================
% APPENDIX F: Interpretability — ACF Scatter Plots
% ============================================================
\section{Interpretability: Imputation vs.\ ACF Change}
\label{app:interpretability}

\begin{figure*}[t]
\centering
\begin{subfigure}[b]{0.78\textwidth}
\includegraphics[width=\textwidth]{figures/fig_acf_aki.png}
\caption{AKI ($\rho{=}0.81$)}
\end{subfigure}\\[0.8em]
\begin{subfigure}[b]{0.78\textwidth}
\includegraphics[width=\textwidth]{figures/fig_acf_sepsis.png}
\caption{Sepsis ($\rho{=}0.74$)}
\end{subfigure}
\caption{Feature imputation fraction vs.\ adaptation-induced autocorrelation change ($\Delta$ACF). Laboratory values (high imputation, upper quadrant) undergo large temporal restructuring; continuously-monitored vitals (lower left) are preserved. Spearman $\rho$ indicates the correlation between imputation fraction and ACF change, both $p < 10^{-9}$.}
\label{fig:acf}
\end{figure*}


% ============================================================
% APPENDIX G: Systematic Method Comparison
% ============================================================
\section{Systematic Method Comparison}
\label{app:comparison}

\Cref{tab:comparison} compares representative methods from five categories across seven dimensions that collectively define the desiderata for clinical time-series domain adaptation with a frozen predictor.
No prior method satisfies all seven.

\begin{table*}[t]
\centering
\caption{Systematic comparison of domain adaptation and frozen-model methods across key desiderata.
\textbf{Model Frozen}: entire target predictor weights fixed (\cmark), partially fixed (P), or trainable (\xmark).
\textbf{Adapt.\ Space}: Input (I), Representation (R), Weight (W), or Output (O) space.
\textbf{Sample-Cond.}: transformation depends on the input instance.
\textbf{Retrieval}: uses a domain-specific datastore for per-instance conditioning.
\textbf{Var.-Len.\ TS}: handles variable-length time series natively.
\textbf{Clinical Tasks}: number of distinct clinical prediction tasks evaluated ($0$ = non-clinical only).
Only our method satisfies all desiderata simultaneously.}
\label{tab:comparison}
\vskip 0.1in
\footnotesize
\setlength{\tabcolsep}{3.2pt}
\begin{tabular}{@{}lcccccc@{}}
\toprule
\textbf{Method} &
\rotatebox{60}{\textbf{Model Frozen}} &
\rotatebox{60}{\textbf{Adapt.\ Space}} &
\rotatebox{60}{\textbf{Sample-Cond.}} &
\rotatebox{60}{\textbf{Retrieval}} &
\rotatebox{60}{\textbf{Var.-Len.\ TS}} &
\rotatebox{60}{\textbf{Clinical Tasks}} \\
\midrule
\multicolumn{7}{@{}l}{\textit{Clinical DA}} \\
VRADA~\citep{purushotham2017vrada}   & \xmark & R & \cmark & \xmark & \cmark & 1 \\
OTTEHR~\citep{ottehr2025}            & \xmark & R & \xmark & \xmark & \xmark & 3 \\
Mutnuri et al.~\citep{mutnuri2024}   & \xmark & W & \cmark & \xmark & \cmark & 3 \\
\midrule
\multicolumn{7}{@{}l}{\textit{General TS-DA}} \\
CoDATS~\citep{wilson2020codats}       & \xmark & R & \cmark & \xmark & \xmark & 0 \\
RAINCOAT~\citep{he2023raincoat}       & \xmark & R & \cmark & \xmark & \xmark & 0 \\
CLUDA~\citep{ozyurt2023cluda}         & \xmark & R & \cmark & \xmark & \xmark & 0 \\
\midrule
\multicolumn{7}{@{}l}{\textit{Frozen-model / Input-space}} \\
VPT~\citep{jia2022vpt}                & \cmark & I & \xmark & \xmark & \xmark & 0 \\
Time-LLM~\citep{jin2024timellm}       & \cmark & I & \xmark & \xmark & \xmark & 0 \\
TATO~\citep{qiu2026tato}              & \cmark & I & P      & \xmark & \cmark & 0 \\
\midrule
\multicolumn{7}{@{}l}{\textit{Retrieval-augmented}} \\
$k$NN-MT~\citep{khandelwal2021knnmt}  & \cmark & O & \cmark & \cmark & \cmark & 0 \\
TS-RAG~\citep{ning2025tsrag}          & \xmark & R & \cmark & \cmark & \cmark & 0 \\
\midrule
\multicolumn{7}{@{}l}{\textit{Foundation model}} \\
ICareFM~\citep{burger2025icarefm}     & \xmark & W & \cmark & \xmark & \cmark & 4 \\
\midrule
\textbf{\retr{} (Ours)} & \cmark & \textbf{I} & \cmark & \cmark & \cmark & \textbf{5} \\
\bottomrule
\end{tabular}
\end{table*}

Representation-space (R) and weight-space (W) methods are architecturally coupled to the model's internals and cannot be formulated without access to intermediate representations or weight tensors, violating the frozen-model constraint.
Input-space methods (VPT, Time-LLM, TATO) keep the model frozen but use input-independent perturbations (VPT, Time-LLM) or handcrafted transforms (TATO), limiting their capacity for heterogeneous clinical data where dozens of features shift simultaneously in different ways.
$k$NN-MT~\citep{khandelwal2021knnmt} is closest in spirit (frozen model, retrieval, per-timestep conditioning) but operates in output space (interpolating token distributions) rather than input space.
\retr{} is the only method combining a fully frozen predictor, learned input-space adaptation, sample-conditional neural transformation, target-domain retrieval, variable-length support, and evaluation across five clinical tasks.


% ============================================================
% APPENDIX H: Computational Cost
% ============================================================
\section{Computational Cost}
\label{app:compute}

\begin{table}[t]
\centering
\caption{Training time per task (single seed, V100S-32GB). Phase~1 checkpoints are reusable across experiments with matching architecture.}
\label{tab:compute}
\small
\begin{tabular}{@{}lrrr@{}}
\toprule
\textbf{Task} & \textbf{Phase 1 (h)} & \textbf{Phase 2 (h)} & \textbf{Total (h)} \\
\midrule
Mortality & 4.5  & 9.0  & 13.5 \\
AKI       & 6.8  & 15.8 & 22.5 \\
Sepsis    & 6.8  & 13.5 & 20.3 \\
LoS       & 7.5  & 15.0 & 22.5 \\
KF        & 2.5  & 5.0  & 7.5 \\
\bottomrule
\end{tabular}
\end{table}

\Cref{tab:compute} reports single-seed training time per task. Phase~1 checkpoints are reusable, amortizing pretraining cost across experiments.

\paragraph{Parameter counts.}
Adapter: 2.59--2.88M parameters depending on task (dominated by 7 AxialBlocks at 265K each, plus 2--3 cross-attention blocks).
Frozen LSTMs: 170K (Sepsis) to 1.42M (AKI).
The adapter-to-model ratio ranges from 2.0$\times$ (AKI) to 15.2$\times$ (Sepsis), but the adapter operates in input space and never modifies the frozen model's parameters.

\paragraph{Inference overhead.}
At test time, the memory bank is pre-built once ($\sim$6 minutes for AKI on V100S).
Per-stay inference latency is $\sim$15.5\,ms (AKI, including $k$-NN retrieval from a 424K-window bank), representing $\sim$9$\times$ overhead over the frozen LSTM alone; this is acceptable for clinical decision support where predictions are computed per hour.

\paragraph{Total compute budget.}
Full development required $\sim$2,400 GPU-hours across $\sim$700 experiments (architecture search, ablation studies, multi-seed validation, DA baselines, \hirid{} experiments).
Minimum reproduction of the best configuration per task (single seed, all five EHR tasks plus five AdaTime datasets): $\sim$107 GPU-hours on V100S-32GB.
Hardware: NVIDIA V100S-32GB (primary), A6000-48GB, RTX 3090-24GB, L40S-48GB, and A100-40GB (SLURM cluster).

\paragraph{AdaTime.}
Per-dataset training: $\sim$21\,h across 10 DA scenarios $\times$ 5 seeds.
Total AdaTime compute: $\sim$108 GPU-hours for full 5-seed validation across all 5 datasets.
